\documentclass[11pt,a4paper]{letter}
\usepackage[T1]{fontenc}
\usepackage[utf8]{inputenc}
\usepackage[french]{babel}
\usepackage{geometry}
\geometry{left=2.5cm, right=2.5cm, top=2cm, bottom=2cm}
\usepackage{hyperref}
\hypersetup{colorlinks=true, urlcolor=blue}
\usepackage{parskip}

\signature{DKHILI Fares}
\address{DKHILI Fares \\ ElManar-Tunis, Tunisie \\ +216 27 491 238 \\ Dkhili.fares2002@gmail.com}

\begin{document}
\begin{letter}{Madame, Monsieur \\ Service Recrutement \\ [Nom de l'entreprise] \\ [Adresse]}

\opening{Madame, Monsieur,}

Dipl\^om\'e en ing\'enierie informatique de la Facult\'e des Sciences de Tunis, avec une sp\'ecialisation en syst\`emes embarqu\'es, je me permets de vous soumettre ma candidature au poste d'\textbf{Ing\'enieur Logiciel Embarqu\'e}.

Au cours de mon stage de fin d'\'etudes chez \textbf{STMicroelectronics}, j'ai men\'e un projet de benchmark compl\`ete de l'\'ecosyst\`eme STM32Cube face \`a 5 plateformes concurrentes (TI, NXP, Renesas, GigaDevice). Ce travail m'a permis de ma\^itriser le d\'eveloppement firmware sur ARM Cortex-M en bare-metal et sous FreeRTOS, le portage de code entre architectures MCU, et l'analyse approfondie de performance et d'empreinte m\'emoire sur des stacks middleware critiques (USB, FatFS, RTOS).

Ce qui me distingue particuli\`erement est ma capacit\'e \`a prendre des initiatives techniques de mani\`ere autonome. Sans demande de mon encadrant, j'ai con\c{c}u et d\'evelopp\'e deux outils d'automatisation : un \textbf{analyseur d'empreinte firmware pilot\'e par l'IA} avec interface chatbot qui automatise l'analyse des \'ecarts de taille et g\'en\`ere des recommandations d'optimisation, et un \textbf{agent Copilot VS Code personnalis\'e} qui automatise la cr\'eation de projets, le portage multi-vendeurs et la configuration de build/debug. Ces outils ont significativement acc\'el\'er\'e le workflow de benchmarking.

Pr\'ealablement, j'ai port\'e le benchmark \textbf{EEMBC CoreMark-PRO} (d\'ependant de POSIX/Linux) vers STM32 en bare-metal --- un d\'efi technique qui a n\'ecessit\'e l'adaptation compl\`ete des couches d'abstraction MITH pour l'architecture ARM.

Mes comp\'etences couvrent : le d\'eveloppement C/C++ embarqu\'e, les RTOS, les pilotes p\'eriph\'eriques (UART, SPI, I2C, CAN), la cross-compilation, Buildroot/Yocto, et les outils STM32 (CubeMX, CubeIDE, IAR). Je suis \'egalement certifi\'e CCNA et form\'e au Deep Learning (NVIDIA).

Motiv\'e par les d\'efis techniques et dot\'e d'une forte capacit\'e d'autonomie, je suis convaincu de pouvoir apporter une contribution significative \`a votre \'equipe de d\'eveloppement embarqu\'e.

Je reste \`a votre disposition pour un entretien \`a votre convenance.

\closing{Cordialement,}

\end{letter}
\end{document}
